% GENERATED FILE. Edit canonical lessons, then run npm run generate:books.
% canonical-source-hash: fnv1a64:cb08d9be5bf6634e
% canonical-lessons: IT-C128-pericoloso, IT-C128-bagnato, IT-C128-asciutto, IT-C128-fresco, IT-C128-morbido

\chapter{Dangerous, Wet, Dry, Cool, Soft}
\label{ch:it-dangerous-wet-dry-cool-soft}
\hlchaptermodality{128}

\begin{chapteropening}
\textbf{By the end of this chapter:} \emph{I can say dangerous, wet, dry, cool and soft.}

Complete the last lesson of chapter 128: I can say dangerous, wet, dry, cool and soft.
\end{chapteropening}

\section[pericoloso]{pericoloso — dangerous}

\label{lesson:IT-C128-pericoloso}



\begin{quote}
Before the new one: say the Italian for ready, then the Italian for safe; sure.
\end{quote}

\subsection*{You'll want to know: pericoloso}
\textbf{pericoloso} — \textquotedblleft{}dangerous\textquotedblright{}.

\begin{cousinweb}
\textbf{pericoloso} $\leftarrow$ Latin \emph{periculosus}, \textquotedblleft{}dangerous\textquotedblright{} (Lewis \& Short) — a learned borrowing; EtymDB 2.1 has no row for \emph{pericoloso}.
\end{cousinweb}

\begin{grammarlens}[title={What you've built}]
One more word for this chapter.
\end{grammarlens}

\subsection*{Your turn}
\begin{itemize}
\raggedright
  \item \emph{Say it:} \emph{pericoloso}
  \item \emph{Say it:} \emph{pericoloso}, once more
  \item \emph{Say it:} say \emph{pericoloso}
  \item \emph{Recall:} say the Italian for ready, then the Italian for safe; sure, then say \emph{pericoloso} again
\end{itemize}

\begin{tcolorbox}[breakable,colback=teal!4,colframe=teal!35!black,title={Before you move on}]
What does pericoloso mean? (\textquotedblleft{}Dangerous\textquotedblright{}.) Say it once more.
\end{tcolorbox}
\section[bagnato]{bagnato — wet}

\label{lesson:IT-C128-bagnato}



\begin{quote}
Before the new one: say the Italian for safe; sure, then the Italian for dangerous.
\end{quote}

\subsection*{You'll want to know: bagnato}
\textbf{bagnato} — \textquotedblleft{}wet\textquotedblright{}.

\begin{cousinweb}
\textbf{bagnato} is the past participle of the verb for \textquotedblleft{}to wet, to bathe\textquotedblright{}, which $\leftarrow$ Vulgar Latin \emph{baneare}, from Late Latin \emph{balneare}, \textquotedblleft{}to bathe\textquotedblright{}. (EtymDB 2.1)
\end{cousinweb}

\begin{grammarlens}[title={What you've built}]
One more word for this chapter.
\end{grammarlens}

\subsection*{Your turn}
\begin{itemize}
\raggedright
  \item \emph{Say it:} \emph{bagnato}
  \item \emph{Say it:} \emph{bagnato}, once more
  \item \emph{Say it:} say \emph{bagnato}
  \item \emph{Recall:} say the Italian for safe; sure, then the Italian for dangerous, then say \emph{bagnato} again
\end{itemize}

\begin{tcolorbox}[breakable,colback=teal!4,colframe=teal!35!black,title={Before you move on}]
What does bagnato mean? (\textquotedblleft{}Wet\textquotedblright{}.) Say it once more.
\end{tcolorbox}
\section[asciutto]{asciutto — dry}

\label{lesson:IT-C128-asciutto}



\begin{quote}
Before the new one: say the Italian for dangerous, then the Italian for wet.
\end{quote}

\subsection*{You'll want to know: asciutto}
\textbf{asciutto} — \textquotedblleft{}dry\textquotedblright{}.

\begin{cousinweb}
\textbf{asciutto} $\leftarrow$ Latin \emph{exsuctus}, \textquotedblleft{}sucked out, dried\textquotedblright{} — inherited. (EtymDB 2.1)
\end{cousinweb}

\begin{grammarlens}[title={What you've built}]
One more word for this chapter.
\end{grammarlens}

\subsection*{Your turn}
\begin{itemize}
\raggedright
  \item \emph{Say it:} \emph{asciutto}
  \item \emph{Say it:} \emph{asciutto}, once more
  \item \emph{Say it:} say \emph{asciutto}
  \item \emph{Recall:} say the Italian for dangerous, then the Italian for wet, then say \emph{asciutto} again
\end{itemize}

\begin{tcolorbox}[breakable,colback=teal!4,colframe=teal!35!black,title={Before you move on}]
What does asciutto mean? (\textquotedblleft{}Dry\textquotedblright{}.) Say it once more.
\end{tcolorbox}
\section[fresco]{fresco — cool; fresh}

\label{lesson:IT-C128-fresco}



\begin{quote}
Before the new one: say the Italian for wet, then the Italian for dry.
\end{quote}

\subsection*{You'll want to know: fresco}
\textbf{fresco} — \textquotedblleft{}cool; fresh\textquotedblright{}.

\begin{cousinweb}
\textbf{fresco} $\leftarrow$ Germanic \emph{frisk-}, \textquotedblleft{}fresh\textquotedblright{} — borrowed, through Medieval Latin \emph{friscus}. (EtymDB 2.1) English \emph{fresh} is its cousin.
\end{cousinweb}

\begin{grammarlens}[title={What you've built}]
One more word for this chapter.
\end{grammarlens}

\subsection*{Your turn}
\begin{itemize}
\raggedright
  \item \emph{Say it:} \emph{fresco}
  \item \emph{Say it:} \emph{fresco}, once more
  \item \emph{Say it:} say \emph{fresco}
  \item \emph{Recall:} say the Italian for wet, then the Italian for dry, then say \emph{fresco} again
\end{itemize}

\begin{tcolorbox}[breakable,colback=teal!4,colframe=teal!35!black,title={Before you move on}]
What does fresco mean? (\textquotedblleft{}Cool; fresh\textquotedblright{}.) Say it once more.
\end{tcolorbox}
\section[morbido]{morbido — soft}

\label{lesson:IT-C128-morbido}



\begin{quote}
Before the new one: say the Italian for dry, then the Italian for cool; fresh.
\end{quote}

\subsection*{You'll want to know: morbido}
\textbf{morbido} — \textquotedblleft{}soft\textquotedblright{}.

\begin{cousinweb}
\textbf{morbido} $\leftarrow$ Latin \emph{morbidus}, \textquotedblleft{}sickly\textquotedblright{}, from \emph{morbus}, \textquotedblleft{}disease\textquotedblright{} — a learned borrowing that softened in meaning. EtymDB 2.1 has no row for \emph{morbido}; it derives French \emph{morbide} and Spanish \emph{mórbido} from \emph{morbidus}.
\end{cousinweb}

\begin{grammarlens}[title={What you've built}]
One more word for this chapter.
\end{grammarlens}

\subsection*{Your turn}
\begin{itemize}
\raggedright
  \item \emph{Say it:} \emph{morbido}
  \item \emph{Say it:} \emph{morbido}, once more
  \item \emph{Say it:} say \emph{morbido}
  \item \emph{Recall:} say the Italian for dry, then the Italian for cool; fresh, then say \emph{morbido} again
\end{itemize}

\begin{tcolorbox}[breakable,colback=teal!4,colframe=teal!35!black,title={Before you move on}]
What does morbido mean? (\textquotedblleft{}Soft\textquotedblright{}.) Say it once more.
\end{tcolorbox}
